\begin{table}[t]
\centering
\small
\setlength{\tabcolsep}{1.2ex}
\begin{tabular}{lccccc}
\toprule
\textbf{智能体} & \textbf{玩家中位数} & \textbf{玩家均值} & \textbf{纪录均值} & \textbf{截断纪录均值} \\
\midrule
DreamerV2 & 1.64 & 11.33 & 0.36 & 0.25 \\
\rlap{No Layer Norm} & 1.66 & \hphantom{0}5.95 & 0.38 & 0.25 \\
\rlap{No Reward Gradients} & 1.68 & \hphantom{0}6.18 & 0.37 & 0.24 \\
\rlap{No Discrete Latents} & 1.08 & \hphantom{0}3.71 & 0.24 & 0.19 \\
\rlap{No KL Balancing} & 0.84 & \hphantom{0}3.49 & 0.19 & 0.16 \\
\rlap{No Policy Reinforce} & 0.69 & \hphantom{0}2.74 & 0.16 & 0.15 \\
\rlap{No Image Gradients} & 0.04 & \hphantom{0}0.31 & 0.01 & 0.01 \\
\bottomrule
\end{tabular}
\caption{DreamerV2 消融实验，按 Atari 200M 帧性能衡量，并按最后一列排序。本实验使用的 DreamerV2 版本比\cref{tab:agg}中的版本略早。每个消融项只移除 DreamerV2 智能体的一部分。离散潜变量和 KL 平衡对 DreamerV2 的成功有重要贡献。此外，世界模型依靠图像梯度学习能带来成功行为的通用表征，即使这些表征并非专门为预测过去奖励而学习。}
\label{tab:abl}
\vspace*{-2ex}
\end{table}
